% Part: sets-functions-relations
% Chapter: sets
% Section: unions-and-intersections
% Hindi translation (hi-Deva-IN) of Open Logic commit 9620cc73f9c8e0ad003c514a5d3748f29611c4c0.

\documentclass[../../../include/open-logic-section]{subfiles}

\begin{document}

\olfileid[hi]{sfr}{set}{uni}
\olsection{सम्मिलन और सर्वनिष्ठ}

\begin{explain}
\olref[sfr][set][bas]{sec} में हमने साझा गुणधर्म द्वारा समुच्चयों की
परिभाषाएँ प्रस्तुत की थीं, अर्थात् $\Setabs{x}{\phi(x)}$ रूप की परिभाषाएँ।
यहाँ हम किसी गुणधर्म~$\phi$ का उपयोग करते हैं, और यह गुणधर्म पहले से
परिभाषित समुच्चयों का उल्लेख कर सकता है। उदाहरण के लिए, यदि $A$ और~$B$
समुच्चय हैं, तो समुच्चय $\Setabs{x}{x \in A \lor x \in B}$ में वे सभी
वस्तुएँ हैं जो $A$ या~$B$ के !!{element} हैं; अर्थात् यह
$A$ और~$B$ के !!{element}ों को एक साथ सम्मिलित करने वाला समुच्चय है। इसे
\olref{fig:union} की तरह देखा जा सकता है, जहाँ रंगीन क्षेत्र दोनों
समुच्चयों $A$ और~$B$ के !!{element}ों को एक साथ दर्शाता है।

\begin{figure}
  \olasset{assets/diagrams/union.tikz}
  \caption{दो समुच्चयों का सम्मिलन $A \cup B$ उन सभी !!{element}ों का
   समुच्चय है जो $A$ या~$B$ में हैं।}
  \ollabel{fig:union} 
\end{figure}

समुच्चयों को इस प्रकार सम्मिलित करने की संक्रिया अत्यंत उपयोगी और सामान्य
है, इसलिए हम इसे एक औपचारिक नाम और संकेत देते हैं।
\end{explain}

\begin{defn}[सम्मिलन]
दो समुच्चयों $A$ और $B$ का \emph{सम्मिलन}, जिसे $A \cup B$ लिखते हैं,
उन सभी वस्तुओं का समुच्चय है जो $A$, $B$, अथवा दोनों के !!{element} हैं।
\[
A \cup B = \Setabs{x}{x \in A \lor x \in B}
\]
\end{defn}

\begin{ex}
क्योंकि किसी !!{element} की पुनरावृत्ति से कोई अंतर नहीं पड़ता, जिन दो
समुच्चयों में कोई !!a{element} उभयनिष्ठ है, उनके सम्मिलन में वह !!{element}
केवल एक बार आता है; उदाहरणतः,
$\{ a, b, c\} \cup \{ a, 0, 1\} = \{a, b, c, 0, 1\}$।

किसी समुच्चय और उसके किसी उपसमुच्चय का सम्मिलन बड़ा समुच्चय ही होता है:
$\{a,
b, c \} \cup \{a \} = \{a, b, c\}$।

किसी समुच्चय का रिक्त समुच्चय के साथ सम्मिलन उसी समुच्चय के समान होता है:
$\{a,
b, c \} \cup \emptyset = \{a, b, c \}$।
\end{ex}

\begin{prob}
सिद्ध कीजिए कि यदि $A \subseteq B$, तो $A \cup B = B$।
\end{prob}

\begin{explain}
हम सम्मिलन की एक ``द्वैत'' संक्रिया पर भी विचार कर सकते हैं। यह संक्रिया
उन सभी !!{element}ों का समुच्चय बनाती है जो~$A$ के !!{element} हैं और
साथ ही~$B$ के भी !!{element} हैं। इस संक्रिया को \emph{सर्वनिष्ठ} कहते
हैं और इसे \olref{fig:intersection} की तरह दर्शाया जा सकता है।
\begin{figure}
  \olasset{assets/diagrams/intersection.tikz}
  \caption{दो समुच्चयों का सर्वनिष्ठ $A \cap B$ उन !!{element}ों का
    समुच्चय है जो दोनों में उभयनिष्ठ हैं।}
  \ollabel{fig:intersection}
\end{figure}
\end{explain}

\begin{defn}[सर्वनिष्ठ]
दो समुच्चयों $A$ और $B$ का \emph{सर्वनिष्ठ}, जिसे $A \cap B$ लिखते हैं,
उन सभी वस्तुओं का समुच्चय है जो $A$ और~$B$ दोनों के !!{element} हैं।
\[
A \cap B = \Setabs{x}{x \in A \land x \in B}
\]
यदि दो समुच्चयों का सर्वनिष्ठ रिक्त हो, तो उन्हें \emph{असंयुक्त} कहते
हैं। इसका अर्थ है कि उनका कोई !!{element} उभयनिष्ठ नहीं है।
\end{defn}

\begin{ex}
यदि दो समुच्चयों का कोई !!{element} उभयनिष्ठ न हो, तो उनका सर्वनिष्ठ
रिक्त होता है: $\{ a, b, c\} \cap \{ 0, 1\} = \emptyset$।

यदि दो समुच्चयों के कुछ !!{element} उभयनिष्ठ हों, तो उनका सर्वनिष्ठ उन
सभी का समुच्चय होता है:
$\{a, b, c \} \cap \{a, b, d \} = \{a, b\}$।

किसी समुच्चय का उसके किसी उपसमुच्चय के साथ सर्वनिष्ठ छोटा समुच्चय ही होता
है: $\{a, b, c\} \cap \{a, b\} = \{a, b\}$।

किसी भी समुच्चय का रिक्त समुच्चय के साथ सर्वनिष्ठ रिक्त होता है:
$\{a, b, c \}
\cap \emptyset = \emptyset$।
\end{ex}

\begin{prob}
विधिपूर्वक सिद्ध कीजिए कि यदि $A \subseteq B$, तो $A \cap B = A$।
\end{prob}

\begin{explain}
हम दो से अधिक समुच्चयों का सम्मिलन या सर्वनिष्ठ भी बना सकते हैं। सामान्य
स्थिति को सँभालने का एक सुगठित तरीका यह है: मान लीजिए कि जिन सभी समुच्चयों
का सम्मिलन (या सर्वनिष्ठ) बनाना है, उन्हें एक ही समुच्चय में एकत्र कर लिया
गया है। तब मूल समुच्चयों के सम्मिलन को उन सभी वस्तुओं का समुच्चय कह सकते
हैं जो इस नए समुच्चय के कम-से-कम एक !!{element} में हैं, और सर्वनिष्ठ को
उन सभी वस्तुओं का समुच्चय कह सकते हैं जो उसके प्रत्येक !!{element} में हैं।
\end{explain}

\begin{defn}
यदि $A$ समुच्चयों का समुच्चय है, तो $\bigcup A$, $A$ के !!{element}ों के
!!{element}ों का समुच्चय है:
\begin{align*}
\bigcup A & = \Setabs{x}{x \text{ निम्न समुच्चय के किसी !!a{element} का अवयव है: } A},
\text{ अर्थात्}\\
& = \Setabs{x}{\text{कोई ऐसा } B \in A
  \text{ है कि } x \in B}
\end{align*}
\end{defn}

\begin{defn}
यदि $A$ समुच्चयों का समुच्चय है, तो $\bigcap A$ उन वस्तुओं का समुच्चय
है जो~$A$ के सभी अवयवों में उभयनिष्ठ हैं:
\begin{align*}
\bigcap A & = \Setabs{x}{x \text{ निम्न समुच्चय के प्रत्येक !!{element} का अवयव है: } A},
\text{ अर्थात्}\\
 & = \Setabs{x}{\text{प्रत्येक } B \in A, x \in B}
\end{align*}
\end{defn}

\begin{ex}
मान लीजिए $A = \{ \{ a, b \}, \{ a, d, e \}, \{ a, d \} \}$।
तब $\bigcup A = \{ a, b, d, e \}$ और $\bigcap A = \{ a \}$।
\end{ex}
\begin{prob}
	दिखाइए कि यदि $A$ समुच्चय है और $A \in B$, तो $A \subseteq \bigcup B$।
\end{prob}

समुच्चयों के अनुक्रम $A_1$, $A_2$, \dots के लिए भी हम यही कर सकते हैं:
\begin{align*}
\bigcup_i A_i & = \Setabs{x}{x \text{ निम्न में से किसी एक समुच्चय का अवयव है: } A_i}\\
\bigcap_i A_i & = \Setabs{x}{x \text{ निम्न प्रत्येक समुच्चय का अवयव है: } A_i}.
\end{align*}

जब हमारे पास समुच्चयों का कोई \emph{सूचकांक} हो, अर्थात् ऐसा समुच्चय $I$
कि हम प्रत्येक $A_i$ पर विचार कर रहे हों जहाँ $i \in I$, तब हम इन
संक्षिप्त संकेतों का भी उपयोग कर सकते हैं:
\begin{align*}
	\bigcup_{i \in I} A_i & = \bigcup \Setabs{A_i }{i \in I}\\
	\bigcap_{i \in I} A_i & = \bigcap\Setabs{A_i}{i \in I}
\end{align*}

अंततः हम~$A$ के उन सभी !!{element}ों के समुच्चय पर विचार कर सकते हैं जो~$B$
में नहीं हैं। इसे \olref{difference} की तरह दर्शाया जा सकता है।

\begin{figure}
  \olasset{assets/diagrams/difference.tikz}
  \caption{दो समुच्चयों का अंतर $A \setminus B$,~$A$ के उन !!{element}ों
    का समुच्चय है जो~$B$ के !!{element} नहीं हैं।}
  \ollabel{difference}
\end{figure}

\begin{defn}[अंतर]
\emph{समुच्चय-अंतर}~$A \setminus B$, $A$ के उन सभी !!{element}ों का समुच्चय
है जो~$B$ के भी !!{element} नहीं हैं, अर्थात्,
\[
A\setminus B = \Setabs{x}{x\in A \text{ और } x \notin B}.
\]
\end{defn}

\begin{prob}
	सिद्ध कीजिए कि यदि $A \subsetneq B$, तो $B \setminus A \neq \emptyset$।
\end{prob}

\end{document}
